\documentclass[12pt, a4paper]{article} % Classe: article, report o book [8]

% --- Lingua e Codifica (Supporto Italiano) ---
\usepackage[utf8]{inputenc}  % Codifica input
\usepackage[T1]{fontenc}      % Codifica font
\usepackage[italian]{babel}   % Lingua italiana [2]

% --- Impaginazione e Margini ---
\usepackage[margin=2.5cm]{geometry} % Imposta margini (2.5cm su tutti i lati)

% --- Pacchetti Scientifici/Matematici ---
\usepackage{amsmath, amssymb, amsfonts, amsthm} % Matematica avanzata [2]
\usepackage{graphicx} % Per inserire immagini
\usepackage{hyperref} % Per collegamenti ipertestuali
\usepackage{booktabs} % Per tabelle professionali

% --- Formattazione Testo e Altro ---
\usepackage[document]{ragged2e}
\usepackage{xcolor}    % Per usare i colori 


\title{Mattia Maltoni}
\date{11-12-2025}

\begin{document}

\maketitle

\section{Ruolo}
Mattia Maltoni è un ingegnere di progetto con esperienza di circa 3 anni
all'interno dell'azienda e si occupa di :
\begin{itemize}
    \item Progettazione e gestione di impianti industriali.
    \item Collaborazione con il reparto commerciale per la definizione delle specifiche
          tecniche.
    \item Supervisione della fase di installazione e messa in servizio degli impianti.
    \item Coordinamento con i fornitori e gestione delle risorse tecniche.
    \item Assistenza tecnica post-vendita e supporto ai clienti.
\end{itemize}

\section{Contatto con l'attuale programma:}
Mattia non utilizza ne consulta Access per le sue attività lavorative. Il
contatto con il prodotto dei commerciali avviene tramite documenti cartacei ed
eventualmente scansionati. Il contratto viene consultato dall'ufficio project
manager per:
\begin{itemize}
    \item Definire il venduto e le specifiche tecniche.
    \item Revisionare con note tecniche il contratto commerciale. Se non presenti
          infomazioni fondamnetali per la progettazione richiedono chiarimenti ai
          commerciali.
    \item Prendere i contatti con i clienti per quanto riguarda le specifiche tecniche e
          le conferme dei preliminari di progetto.
\end{itemize}

\section{Gestione temporale e monitoraggio}
Il sistema deve permettere un controllo granulare delle date per identificare
colli di bottiglia e garantire il rispetto delle consegne.

\begin{itemize}
    \item \textbf{Allineamento date emissione/stampa:}
          Il sistema deve registrare e confrontare:
          \begin{itemize}
              \item \textit{Data emissione ordine:} Data di conferma da parte del cliente.
              \item \textit{Data stampa/presa in carico:} Data di effettiva ricezione interna.
          \end{itemize}
          Il software deve evidenziare automaticamente eventuali discrepanze temporali (ritardi) tra queste due date.

    \item \textbf{Pianificazione a ritroso (Back-Scheduling):}
          Partendo dalla \textit{Data evasione} (spedizione prevista) o \textit{Start-up impianto}, il sistema deve calcolare le milestone intermedie per assicurare che tutte le parti siano processate nei tempi corretti.
\end{itemize}

\section{Digitalizzazione}
Eliminazione del supporto cartaceo e delle annotazioni manuali non tracciabili.

\begin{itemize}
    \item \textbf{Transizione digitale:} Tutti i campi, le note e le variazioni devono essere compilabili esclusivamente via software.
    \item \textbf{Log delle modifiche:} Ogni variazione ai dati di commessa deve essere storicizzata. Il sistema deve rispondere ai requisiti di tracciabilità:
          \begin{itemize}
              \item \textit{Chi} ha effettuato la modifica.
              \item \textit{Quando} è stata effettuata (timestamp).
              \item \textit{Cosa} è stato modificato (valore precedente vs valore attuale).
          \end{itemize}
    \item \textbf{Accessibilità orizzontale:} I dati aggiornati devono essere visibili in tempo reale a tutti gli utenti con i relativi permessi, eliminando le versioni obsolete stampate su carta.
\end{itemize}

\section{Workflow, Stati e Notifiche}
Automazione del flusso di lavoro tra i reparti (Commerciale, Impianti, UT ,
UTE).

\begin{itemize} 
    \item \textbf{Gestione degli stati:} Il software deve distinguere chiaramente lo stato di avanzamento, ad esempio:
          \begin{enumerate}
              \item \textit{Pre-inserimento} (Bozza/Commerciale).
              \item \textit{Inserito in Impianti} (Validato tecnicamente).
          \end{enumerate}
    \item \textbf{Sistema di Alerting:} Al passaggio di stato in "Inserito in Impianti", il sistema deve inviare una notifica automatica (email/alert in-app) ai reparti a valle (es. Reparto Elettrico/UTIE) per segnalare la disponibilità della commessa, sostituendo il passaggio fisico del fascicolo.
\end{itemize}

\section{Anagrafica Compilatore}
\begin{itemize}
    \item \textbf{Attribuzione responsabilità:} Deve essere presente un campo "Compilato da" autocompilato con l'utente loggato.
    \item \textbf{Flessibilità:} Il campo deve rimanere editabile per gestire casi di delega o correzione dell'assegnatario.
\end{itemize}

\section{Verifica Tecnica (Contratto vs Schema)}
Modulo dedicato alla gestione delle discrepanze tra il venduto e l'eseguito
tecnico.

\begin{itemize}
    \item \textbf{Log variazioni tecniche (Contratto tecnico):}
          Il software deve includere una sezione specifica (sostitutiva di file Word/note esterne) per registrare le differenze tra il contratto commerciale e lo schema P\&ID (es. numero valvole, componentistica extra).
    \item \textbf{Obiettivo di integrazione:}
          Creare un flusso dove le modifiche allo schema tecnico aggiornino le specifiche di commessa, mantenendo allineati i dati economici e tecnici.
\end{itemize}

\end{document}